This chapter outlines the framework for our large-scale empirical analysis, which applies the newly developed tool to the complete dataset of \textbf{602} mined software projects. We detail the descriptive measures used to summarise the detector outputs for the seven evaluated SQL antipatterns.

\section{Quantitative Analysis}\label{sec:projectAnalysisQuantitative}

For each of the seven targeted antipatterns, we calculated the following statistical values:

\begin{itemize}
    \item the total number of detected antipattern occurrences,
    \item the number of projects containing that antipattern,
    \item the average number of detected occurrences per project among all \num{602} projects,
    \item the average number of detected occurrences per project among projects containing that antipattern.
\end{itemize}

These quantities describe detector flags rather than verified antipattern prevalence. They are not normalised by repository size or exposure to relevant jOOQ APIs. For antipatterns studied by other works, we therefore compare definitions and qualitative patterns but do not compare rates directly.

The replication repository retains unadjusted exploratory co-flagging matrices. They omit joint-zero units and do not control for repository size, API exposure, or class-specific detector error. We therefore do not use them to answer a research question or infer clusters, dependencies, or predictors.

Based on these results, we provide an answer to \textbf{RQ4} in Section~\ref{sec:resultsRQ4}.

For two of the detected query antipatterns (“Implicit Columns” and “Poor Man's Search Engine”), we analysed the jOOQ API methods associated with each detector flag. We employed an iterative pattern-matching approach, assigning one post hoc manifestation category to each flag by comparing the containing code snippet against a collection of patterns.

We skipped the analysis for all database design antipatterns, since for them, the associations only display how jOOQ generates classes based on an underlying problematic schema, and do not provide insights into potential API design issues.

We also skipped the analysis for the “Fear of the Unknown” antipattern, which we detect from both classes representing the database schema and source classes containing DQL and DML statements. We initially planned to constrain the labelling to only occurrences found in DQL and DML statements, but only a negligible number of “Fear of the Unknown” occurrences (\num{16} occurrences, \qty{1.3}{\percent} of total) were found in DQL and DML statements, not providing a sufficient sample size.

Based on this analysis, we provide an answer to \textbf{RQ5} in Section~\ref{sec:resultsRQ5}. The categories cover \num{6860} of \num{7289} Implicit Columns flags (\qty{94.1}{\percent}) and \num{546} of \num{583} Poor Man's Search Engine flags (\qty{93.7}{\percent}); the remaining flags are reported as other or uncategorised.

\section{Qualitative Analysis}

We discussed the results of quantitative analysis from a qualitative perspective in Sections~\ref{sec:comparison} to~\ref{sec:apiDesign}. We assessed differences between our detector-output distributions and findings from prior works focusing on plain SQL, while accounting for the incompatible units of analysis.

Additionally, we examined the jOOQ API manifestations most commonly assigned to detector flags.
